\worldchapter{Grundregeln}{game-rules}
\label{ch:game-rules}

\begin{multicols*}{2}

Obwohl die Spielleiterin die Bösen steuert, steht sie eigentlich auf der Seite der Helden.
Spieler und Spielleiterin begeben sich zusammen auf eine fantastische Reise, bei der die Erzählung im Vordergrund steht.

Die Spieler steuern ihre \textbf{Charaktere} in dieser fiktiven Welt, und handeln in ihrem Namen.
Die Fähigkeiten eines Charakters, ihr Hab und Gut, ihre Stärken und Schwächen werden auf dem \textbf{Charakterbogen} niedergeschrieben.
Du als Spieler beschreibst der Spielleiterin das Handeln deines Charakters.
Möglicherweise erfordert die Handlung deines Charakters eine \textbf{Probe}, einen Test auf ein bestimmtes \textbf{Attribut}, bei dem passende \textbf{Fertigkeiten} zusätzliche Vorteile geben können.

Diese Proben sind der Kern des Spiels.
Es ist die Grundmechanik, um über den Erfolg oder das Misslingen einer Handlung zu bestimmen.

\section{Die Probe}\label{sec:die-probe}

Die Probe\index{Probe} ist der zentrale Bestandteil der Regeln.
Wenn eine Situation einen ungewissen Ausgang hat, wird eine Probe verwendet, um den Ausgang der Situation zu bestimmen.
Tirakans Reiche benutzt für alle Proben einen hundertseitigen Würfel, einen \textbf{W100}.
In der Regel werden zwei zehnseitige Würfel (\textbf{W10}) geworfen, wobei einer der Würfel die Zehnerstelle und der andere die Einerstelle des Wurfes darstellen.
Ein Ergebnis von \textbf{00} bedeutet 100.

Ob Erkundung, soziale Spannung, Verfolgung, Ritual oder Kampf: Wenn der Ausgang unsicher ist und Konsequenzen hat, wird eine Probe abgelegt.

Die Regeln legen besonderen Wert auf den Ausgang der Probe. Es ist nicht einfach nur ein geschafft oder nicht geschafft, es wird besonderen Wert auf die \textbf{Konsequenzen} einer Probe gelegt. Ob im alltäglichen, im Handwerk oder beim Ausüben von Magie, jede Probe hat Folgen für die Situation oder den Charakter.

\subsection{Zielwert}\label{subsec:zielwert}

Eine Probe wird immer auf genau ein \textbf{Attribut}\index{Attribut} geworfen.
Die Spielleiterin bestimmt das passende Attribut für die Situation, wenn die Regeln nicht etwas anderes vorsehen.
Zu den Attributen zählen \textbf{Geist}, \textbf{Wille}, \textbf{Instinkt}, \textbf{Geschick}, \textbf{Körper}, \textbf{Erscheinung}, \textbf{Gabe} und \textbf{Wahrnehmung}.


Dieser \textbf{Zielwert}\index{Zielwert} muss unterschritten werden, damit die Probe erfolgreich ist. Dieser Zielwert ist gestaffelt nach dem Rang, den der Charakter in dem Attribut hat:\\

\begin{tabular}{ll}
    \toprule
    \textbf{Attribut}  & \textbf{Zielwert}\\
    \midrule
    \textbf{0}        & 30 \\
    \textbf{1}        & 45 \\
    \textbf{2}        & 60 \\
    \textbf{3}        & 75 \\
    \textbf{4}        & 90 \\
    \bottomrule
\end{tabular}

Der Zielwert wird immer auf 5 bis 95 begrenzt.
Eine Probe gelingt, wenn der Wurf kleiner als der Zielwert ist.

\subsection{Fertigkeit}\label{subsec:probe-fertigkeit}

Besitzt der Charakter eine passende Fertigkeit, die er in der Situation einsetzen kann, so gelten die Vorteile der Fertigkeit. Fertigkeiten sind ausführlich in \cref{sec:fertigkeiten} beschrieben. Fertigkeiten entscheiden, ob eine Handlung fachgerecht möglich ist und welche Rangvorteile gelten, verändern den Zielwert aber nicht.

\begin{tabularx}{\linewidth}{lX}
    \toprule
    \textbf{Fertigkeit}  & \textbf{Effekt}\\
    \midrule
    \textbf{0 Improvisiert}        & Bei Misslingen der Probe wird die Konsequenzstufe um eins erhöht.\\
    \textbf{1 Geübt}        &  Knapp gescheiterte Proben zählen als Stillstand und nicht als geringe Konsequenz.\\
    \textbf{2 Erfahren}        & Einmal pro Szene darf eine Probe neu gewürfelt werden, das neue Ergebnis zählt. \\
    \textbf{3 Meisterlich}        & Wenn die Probe gelingt, wird die Ergebnisstufe um eins angehoben.  \\
    \textbf{4 Legendär}        & Einmal pro Szene darf das Ergebnis einer Probe vollständig ignoriert werden. \\
    \bottomrule
\end{tabularx}

Die Stufen der Fertigkeiten gelten kumulativ. Das bedeutet, dass ein Charakter mit einem Rang von 3 in einer Fertigkeit bei einem Wurf auch die Boni der Ränge 1 und 2 bekommt.

\subsection{Schwierigkeit}

Eine Probe kann erschwert sein, wenn sie unter schwierigen Bedingungen ausgeführt wird.
Die Spielleiterin legt hier einen Malus fest, der vor dem Wurf auf den Zielwert angerechnet wird, bevor geworfen wird.

\begin{itemize}[leftmargin=*]
\item \textbf{Leicht}: +20
\item \textbf{Standard}: +0
\item \textbf{Schwer}: -20
\item \textbf{Extrem}: -40
\end{itemize}

\subsection{Ablauf einer Probe}

\begin{enumerate}[leftmargin=*]
\item Der Spieler beschreibt Absicht und Vorgehen.
\item Die Spielleiterin bestimmt das passende Attribut und legt fest, ob eine passende Fertigkeit die Handlung stützt, einschränkt oder improvisiert werden muss.
\item Alle Modifikatoren werden zusammengerechnet.
\item Es wird 1W100 gewürfelt.
\item Ergebnisstufe, unmittelbare Wirkung und Konsequenz werden sofort abgehandelt.
\end{enumerate}

\subsection{Ergebnisstufe}

Das Resultat der Probe ergibt die Ergebnisstufe.
Ist der geworfene Wert bis zu 10 Punkte vom Zielwert entfernt, so ist die Probe knapp gelungen oder misslungen.
Weicht er mehr als 30 Punkte ab, so handelt es sich um einen starken Erfolg oder schweren Fehlschlag.
Ein Wurf von 01--05 ist immer ein kritischer Erfolg, alles von 96--100 ist immer ein katastrophales Scheitern.

\begin{tabular}{ll}
    \toprule
    \textbf{Ergebnisstufe}            & \textbf{Wurf}          \\
    \midrule
    \textbf{Kritischer Erfolg}        & 01--05   \\
    \textbf{Starker Erfolg}           & Marge 30+        \\
    \textbf{Sauberer Erfolg}          & Marge 10--29     \\
    \textbf{Knapp geschafft}          & Marge 0--9       \\
    \textbf{Knapp gescheitert}        & Marge 1--9   \\
    \textbf{Klar gescheitert}         & Marge 10--29 \\
    \textbf{Schwer gescheitert}       & Marge 30+    \\
    \textbf{Katastrophal gescheitert} & 96--00    \\
    \bottomrule
\end{tabular}

\subsection{Konsequenzen}\label{sec:konsequenzen}

Jede misslungene Probe kann zu \textbf{Konsequenzen}\index{Konsequenz} führen, die über die eigentliche Intention der Probe hinaus gehen.
Konsequenzen werden in drei Kategorien eingeteilt und haben unterschiedliche Folgen.
Diese werden von der Spielleiterin festgelegt.

\begin{itemize}[leftmargin=*]
\item \textbf{Gering}: kurze Verzögerung, kleiner Verlust, kurze Irritation.
\item \textbf{Ernst}: Wunde, deutlicher Verlust, Eskalation, schlechtere Position.
\item \textbf{Kritisch}: schwere Wunde, Zielverschiebung, massiver Kollateralschaden, bleibende Folge.
\end{itemize}

In der Regel führt eine bestimmte Ergebnisstufe immer zu einer bestimmten Klasse an Konsequenzen, dies kann natürlich je nach Situation von der Spielleiterin verändert werden.

\begin{itemize}[leftmargin=*]
\item Knapp gescheitert \textrightarrow{} Gering
\item Klar gescheitert \textrightarrow{} Ernst
\item Schwer oder katastrophal gescheitert \textrightarrow{} Kritisch
\item Knapp geschafft in einer verzweifelten Lage \textrightarrow{} Erfolg mit geringer Konsequenz
\end{itemize}

\subsection{Gegenproben}\label{sec:gegenproben}

Wenn zwei Seiten direkt gegeneinander handeln, würfeln beide gegen ihren eigenen Zielwert.

\begin{itemize}[leftmargin=*]
\item Eine Seite gelingt, die andere scheitert: Die erfolgreiche Seite gewinnt.
\item Beide gelingen: Die höhere Erfolgs-Marge gewinnt.
\item Beide scheitern mit 10+ Differenz: die geringere Fehlmarge bekommt die \textbf{Oberhand}\index{Oberhand}.
\item Beide scheitern mit 0--9 Differenz: Stillstand, und die Situation verschlechtert sich fiktional.
\end{itemize}

Oberhand erlaubt genau eines:

\begin{itemize}[leftmargin=*]
\item Initiative behalten.
\item Eine Bewegung um ein Distanzsegment erzwingen.
\item Eine geringe Konsequenz verursachen.
\end{itemize}

\subsection{Reihenfolge bei Gleichstand}

Wenn die Abfolge relevant wird, gilt:

\begin{enumerate}[leftmargin=*]
\item höhere Ergebnisstufe
\item höhere Erfolgs-Marge
\item höherer \textbf{Instinkt}
\item höheres \textbf{Geschick}
\item niedrigerer Rohwurf
\item sonst gleichzeitig
\end{enumerate}

\section{Attribute}\label{sec:attribute}\index{Attribut}

Attribute beschreiben die grundlegenden Anlagen eines Charakters.
Sie geben an, wie zuverlässig jemand allein durch Veranlagung, Körper, Geist oder innere Haltung in einem Bereich handeln kann.

Jede Probe wird auf genau ein Attribut abgelegt.
Dieses Attribut liefert den Zielwert der Probe.
Eine passende Fertigkeit verändert den Zielwert nicht, sondern bestimmt, wie kontrolliert, sicher oder professionell der Charakter die Handlung ausführt.

\subsection{Attributsstufen}

Attribute liegen in der Regel zwischen 0 und 4.
Neue Charaktere beginnen mit Werten von 0 bis 3; durch Fortschritt kann ein Attribut auf 4 steigen.
Jede Attributsstufe hat einen festen Zielwert.

\paragraph{0 -- Schwach}
Der Charakter ist in diesem Bereich deutlich eingeschränkt, ungeübt oder unterdurchschnittlich. Zielwert \textbf{30}.

\paragraph{1 -- Solide}
Der Charakter ist alltagstauglich und verlässlich, aber nicht besonders herausragend. Zielwert \textbf{45}.

\paragraph{2 -- Geübt}
Der Charakter ist klar kompetent und kann sich in belastenden Situationen behaupten. Zielwert \textbf{60}.

\paragraph{3 -- Herausragend}
Der Charakter gehört zu den Besten, die man in seinem Umfeld gewöhnlich antrifft. Zielwert \textbf{75}.

\paragraph{4 -- Außergewöhnlich}
Der Charakter erreicht in diesem Bereich ein seltenes, fast schon legendäres Niveau. Zielwert \textbf{90}.

\subsection{Attribute im Spiel}

Die Spielleiterin wählt immer das eine Attribut, das für Absicht und Vorgehen am wichtigsten ist.
\textbf{Körper} deckt rohe Kraft, Zähigkeit und körperliches Durchhalten ab.
\textbf{Geschick} steht für Präzision, Schnelligkeit und feine Kontrolle.
\textbf{Instinkt} regelt Reaktionen, unmittelbares Gespür und den Umgang mit plötzlicher Gefahr.
\textbf{Geist} beschreibt Analyse, Erinnerung und Verständnis.
\textbf{Wille} trägt Disziplin, Standhaftigkeit und innere Kontrolle.
\textbf{Erscheinung} steht für Präsenz, Druck und soziale Wirkung.
\textbf{Gabe} ist das Attribut für Magie.
\textbf{Wahrnehmung} steht für scharfe Sinne, Aufmerksamkeit für Details und das rechtzeitige Bemerken von Veränderungen.

\begin{pexample}
Jamie will im Regen den Fremden an der Tür genauer betrachten.
Die Spielleiterin verlangt eine Probe auf \textbf{Wahrnehmung}.
Mit Wahrnehmung 2 hat Jamie einen Zielwert von 60.
Besitzt er zusätzlich eine passende Fertigkeit, gelten deren Rangvorteile, ohne dass sich der Zielwert selbst erhöht.
\end{pexample}

\section{Fertigkeiten}\label{sec:fertigkeiten}\index{Fertigkeit}

Fertigkeiten beschreiben Ausbildung, Erfahrung und praktische Kompetenz eines Charakters. Sie ersetzen keine Attribute, sondern zeigen, wie sicher und kontrolliert ein Charakter in einem bestimmten Bereich handeln kann.

Eine Probe wird weiterhin auf das passende Attribut abgelegt. Eine passende Fertigkeit verändert den Zielwert nicht, sondern den Umgang mit Erfolg, Fehlschlag und Konsequenzen.

Hat ein Charakter keine passende Fertigkeit, kann er einfache Handlungen trotzdem versuchen. Solche Proben gelten als improvisiert. Die Spielleiterin kann komplexe, gefährliche oder spezialisierte Anwendungen ausschließen oder mit zusätzlichen Konsequenzen verbinden. Falls die Probe misslingt, steigt die Konsquenzstufe um 1.

\subsection{Fertigkeitsränge}

Die Vorteile der Fertigkeitsränge bauen aufeinander auf. Ein Charakter mit Fertigkeit Rang 3 erhält also auch die Vorteile von Rang 1 und Rang 2.

\paragraph{Rang 0 -- Improvisiert}
Der Charakter hat keine nennenswerte Ausbildung in diesem Bereich. Einfache Anwendungen sind möglich, komplexe oder spezialisierte Anwendungen können unmöglich sein oder zusätzliche Konsequenzen auslösen.

\paragraph{Rang 1 -- Geübt}
Der Charakter kennt die Grundlagen. Knappe Fehlschläge mit dieser Fertigkeit zählen als Stillstand statt als Fehlschlag mit Konsequenz. Die Handlung gelingt nicht, aber die Lage verschlechtert sich nicht wesentlich.

\paragraph{Rang 2 -- Erfahren}
Der Charakter handelt routiniert. Einmal pro Szene darf eine passende Probe neu gewürfelt werden. Das neue Ergebnis gilt.

\paragraph{Rang 3 -- Meisterlich}
Der Charakter arbeitet mit außergewöhnlicher Kontrolle. Wenn eine passende Probe gelingt, wird die Ergebnisstufe um 1 verbessert.

\paragraph{Rang 4 -- Legendär}
Der Charakter zählt zu den wenigen herausragenden Könnern seines Fachs. Einmal pro Szene darf er eine Konsequenz aus einer passenden Probe vollständig ignorieren.

\subsection{Fehlende Fertigkeiten}

Fehlende Fertigkeiten verhindern eine Probe nicht automatisch. Ein Charakter ohne Schwimmen kann versuchen, sich über Wasser zu halten; ein Charakter ohne Heilkunde kann versuchen, eine Blutung zu stoppen.

Die Spielleiterin entscheidet jedoch, ob die Handlung ohne Ausbildung möglich ist. Je spezieller, gefährlicher oder professioneller die Handlung ist, desto eher benötigt sie eine passende Fertigkeit.

\begin{pexample}
Ein Charakter mit Körper 2 und Schwimmen 0 kann versuchen, einen ruhigen Fluss zu durchqueren. Bei starkem Wellengang, schwerer Rüstung oder dem Versuch, eine andere Person zu retten, kann die Spielleiterin die Handlung erschweren oder ohne Schwimmen ausschließen.

Ein Charakter mit Schwimmen 2 darf in derselben Situation einmal pro Szene neu würfeln. Ein Charakter mit Schwimmen 3 verbessert bei Erfolg zusätzlich die Ergebnisstufe. Ein Charakter mit Schwimmen 4 kann einmal pro Szene eine Konsequenz, etwa verlorene Ausrüstung oder Erschöpfung, vollständig vermeiden.
\end{pexample}

\section{Omen}

\textbf{Omen}\index{Omen} ist eine gemeinsame Ressource der Gruppe.
Das Maximum an Omen der Gruppe entspricht 2 plus der Hälfte des Glaubensniveaus (abgerundet).
Zu Sitzungsbeginn wird Omen auf sein Maximum aufgefüllt.

Ein Charakter oder die Gruppe kann 1 Omen ausgeben für:

\begin{itemize}[leftmargin=*]
\item +20 auf den Zielwert vor einem Wurf
\item Neu würfeln und das neue Ergebnis behalten
\item Eine Konsequenz um eine Stufe senken
\item Einmal die Reihenfolge unterbrechen und sofort handeln
\item Eine Szenenfrage stellen und eine klare Wahrheit erhalten
\end{itemize}

Die Szenenfrage darf gestellt werden, wenn der Charakter die Antwort plausibel bemerken, erinnern, deuten oder aus der Lage lesen konnte.
Die Spielleiterin entscheidet, ob die Antwort direkt, eng oder symbolisch ausfällt, aber sie muss wahr sein.

+1 Omen wird gewonnen, wenn eine große Komplikation aus Schwäche, Eid, Schuld oder Verwundbarkeit akzeptiert wird oder wenn eine Bindung mit wirklichem Preis ausgetragen wird.
\end{multicols*}
